\section{Motivating Example: Degraded Fault Distance in Cultivation}
\label{sec:fault_distance_of_magic_state_cultivation}
The reason for the widening T--S discrepancy is that for S-state cultivation,
the double checks have the fault distance they were designed for.
For T-state cultivation,
the final distance-3 (-5) double check
instead has fault distance 2 (3);
this means that the logical error rate per kept shot
scales as $p^2$ ($p^3$) instead of $p^3$ ($p^5$), for noise level $p$.
For the rest of this paper we thus focus on the final double check.
The following shows an explicit example of
two error events that can lead to a logical error
in distance-3 T-state cultivation.

\begin{eg}\label{eg:d3_color_code_acceptance_probability}
Consider two 1-qubit error events, each occurring with probability $\upTheta(p)$;
namely, an $\X_1$ and a $\Y_3$ error event
at the locations shown in \cref{fig:2_fault_configuration}.
The $\X_1$ is a hook error:
it propagates through two CXs to become $\X_{0 \a 1}$,
but ancilla qubit $\a$ is measured in the X-basis
so the X error on it will have no effect.
This leaves us with $\X_{0 1} \Y_3$ immediately before the T-gates,
which will go undetected past this circuit.
Note that the data qubits $(0, 1, 3)$ support a logical Pauli;
we will omit their subscripts in what follows.
Pushing any unitary $U$ through a gate $G$ transforms $U$ into $GUG^\dag$;
now consider pushing $\X\X\Y$ through the given T-gates:
\begin{align}
	E
	&:=
	(\T \T^\dag \T)
	(\X \X \Y)
	(\T \T^\dag \T)^\dag \notag \\
	&=-(\H_+)_0 (\H_-)_{1 3}.
	\label{eq:clifford_error_in_superposition}
\end{align}
The theory we develop in the next section,
together with \cref{lem:insert_projector},
shows that
the Clifford error $E$ flips no subsequent detectors
with probability of 1/4 and,
if so,
leads to an undetected logical error.
An intuitive explanation is:
$E$ completely randomises two X-type stabiliser generators of the code,
so the probability of both returning to their +1 eigenvalue is $(1/2)^2$.
As for the logical information,
$E$ anticommutes with the stabiliser $\overline{H_+}$
that defines the logical state,
so it rotates the logical state to the orthogonal one.
\end{eg}

To contrast,
we briefly show why the error $\X_{0 1} \Y_3$
never leads to logical error in the analogous circuit for S-state cultivation.
\begin{eg}
\label{eg:d3_color_code_s_state}
In S-state cultivation,
the T-gates are replaced by S-gates.
Consider pushing $\X\X\Y$ through these S-gates:
\begin{equation}
	(\S\S^\dag \S) (\X\X\Y)(\S\S^\dag \S)^\dag
	=\Y_{0 1} \X_3.
\end{equation}
This Pauli error anticommutes with $G_2$
from the set of stabiliser generators
\cref{eq:d3_color_stabilizer_generators};
hence, it is deterministically detected.
\end{eg}

The next section is motivated by
the question posed at the end of
\cref{eg:d3_color_code_acceptance_probability}
(\textit{how often does the Clifford error evade detection
and end up as an accepted logical error?});
unfortunately however,
the final double check is followed immediately by the escape stage,
whose details are quite involved.
We avoid propagating the Clifford error $E$ through this stage
using our lemma in \cref{sec:reducing_code_conversion_to_codespace_projection},
which allows us to insert
an imaginary noiseless stabiliser measurement round beforehand
without affecting the post-escape state.
We can then analyse the effect of $E$ on that imaginary round instead,
which is the topic of the next section.
Then, in \cref{sec:thorough_analysis},
we resume our analysis of MSC
by enumerating the low-weight logical errors
in MSC-3 and -5.